% New Section 5. Paragraph roles and source IDs are editorial annotations.
\section{结果}
\label{sec:results}

% Role: scope, links to setup. Evidence: E1w/E1a/E2/E4/E8; frozen protocols.
本节首先检验MCA之后的残差补偿效果，再比较学习方法之间的优势与代价。测试结果沿用同一44名历史工程测试被试，验证集组件比较及方向敏感性分别报告。以下区间均来自既有被试级bootstrap，未校正多重比较，不将均值排序、区间含零或验证集趋势解释为全面优越、统计等效或独立测试确认。

% Role: principal result versus MCA. Evidence: E1w/E2. Relative percentages are arithmetic on means.
\subsection{相对MCA的残差补偿效果}
Hybrid降低了MCA的三项主要频谱误差。表\ref{tab:primary-results}中，全域ERB从1.0822降至0.8175\,dB，对侧$25^\circ$ ERB从1.7458降至1.2004\,dB，对侧高频误差从4.6986降至3.4549\,dB。以MCA均值为分母，三项相对降幅分别为24.46\%、31.24\%和26.47\%；这些百分比描述dB误差指标均值的相对变化，不表示线性声压幅度的下降比例。

% E1a: existing 44-test-subject mean and sample SD; no new aggregation.
\begin{table}[htbp]
\centering\small
\setlength{\tabcolsep}{3pt}
\bicaption[tab:primary-results]{表}{Q26十方法主要指标（44名测试被试，767个插值方向）}{Tab.}{Primary metrics with Q26 (44 test subjects, 767 interpolation directions)}
\begin{tabular}{lrrrr}
\toprule
方法 & 全域ERB & 对侧$25^\circ$ ERB & 对侧高频 & 水平面ILD\\
 & dB$\downarrow$ & dB$\downarrow$ & dB$\downarrow$ & dB$\downarrow$\\
\midrule
SH only & $2.6854 \pm 0.1072$ & $3.8353 \pm 0.4484$ & $9.0669 \pm 0.8240$ & $3.5511 \pm 0.6637$ \\
SUpDEq SH & $1.8821 \pm 0.5060$ & $2.2281 \pm 0.1947$ & $5.9936 \pm 0.3439$ & $2.0181 \pm 1.5643$ \\
SUpDEq NN & $1.8518 \pm 0.2901$ & $2.2405 \pm 0.2123$ & $5.5954 \pm 0.2581$ & $1.6370 \pm 0.4687$ \\
SUpDEq Bary. & $1.7522 \pm 0.2567$ & $2.1821 \pm 0.1929$ & $5.4756 \pm 0.2171$ & $1.6151 \pm 0.4424$ \\
MCA & $1.0822 \pm 0.0959$ & $1.7458 \pm 0.1551$ & $4.6986 \pm 0.2661$ & $0.8294 \pm 0.1582$ \\
MCAR v3.5.1 & $0.8179 \pm 0.1755$ & $1.2745 \pm 0.1477$ & $3.5085 \pm 0.2879$ & $\mathbf{0.6447} \pm 0.2479$ \\
FSP-AE & $1.1845 \pm 0.3492$ & $1.9089 \pm 0.7008$ & $\mathbf{3.1644} \pm 0.6854$ & $0.7587 \pm 1.0397$ \\
RANF & $1.0632 \pm 0.1264$ & $1.5573 \pm 0.1636$ & $3.4697 \pm 0.2538$ & $0.7755 \pm 0.1888$ \\
Hybrid & $0.8175 \pm 0.3427$ & $\mathbf{1.2004} \pm 0.0952$ & $3.4549 \pm 0.2575$ & $0.7794 \pm 0.8389$ \\
Bounded E25 & $\mathbf{0.7941} \pm 0.2397$ & $1.2015 \pm 0.1138$ & $3.4766 \pm 0.2788$ & $0.6744 \pm 0.4881$ \\
\bottomrule
\end{tabular}
\par\smallskip
\parbox{\linewidth}{\footnotesize 数值为被试均值$\pm$样本标准差。加粗仅表示该列最低均值，不表示对其他方法的显著优势。SUpDEq Bary.为重心插值；Hybrid采用E190三个末点成员。}
\end{table}


% Role: paired support and boundary. Evidence: E2, candidate HYBRID/baseline MCA.
三项频谱配对差的95\%区间依次为$[-0.315895,-0.178416]$、$[-0.594202,-0.491940]$和$[-1.297170,-1.187683]$\,dB，均低于零；对应改善人数为43、43和44人。水平面ILD均值也由0.8294降至0.7794\,dB，但差值为$-0.049935$\,dB、区间为$[-0.220991,+0.235298]$\,dB，仍包含零。因此，相对MCA的明确证据集中在频谱重建，不能同时宣称已获得稳定的ILD收益。

% Role: trade-off against the project's strong learned baseline. Evidence: E1a/E2.
\subsection{与其他学习方法的比较}
相对MCAR v3.5.1，Hybrid的优势和代价并存。对侧$25^\circ$ ERB及对侧高频均值差分别为$-0.074093$和$-0.053570$\,dB，95\%区间分别为$[-0.111932,-0.045892]$和$[-0.089022,-0.022387]$\,dB。全域ERB差仅$-0.000393$\,dB且区间含零。与此同时，水平面ILD增加$0.134741$\,dB，区间为$[+0.012705,+0.344849]$\,dB；频谱改善伴随双耳能量关系上的代价。

% Role: strongest competing metrics, not uniform superiority. Evidence: E1a/E2.
十方法的四项主要最低均值分别属于Bounded、Hybrid、FSP-AE和MCAR。Hybrid相对Bounded的全域ERB和水平面ILD差分别为$+0.023413$、$+0.105068$\,dB，配对区间均高于零。其对侧$25^\circ$ ERB均值虽略低，但差值$-0.001121$\,dB的区间$[-0.022054,+0.014415]$\,dB含零；对侧高频则有小幅降低，差值$-0.021669$\,dB、区间$[-0.050092,-0.000766]$\,dB。最低均值因而不等于对所有对照均有清晰区分。

% Role: external learned baselines. Evidence: E2; avoid assigning implementation gains to literature.
Hybrid相对FSP-AE和RANF的全域及对侧$25^\circ$ ERB配对区间均低于零，但FSP-AE的对侧高频误差更小：Hybrid高出$0.290557$\,dB，区间为$[+0.089797,+0.413038]$\,dB。相对RANF的高频差值为$-0.014772$\,dB，区间含零；相对这两个方法的水平面ILD区间也均含零。故本文将Hybrid定位为具有特定频谱优势的残差细化方案，而非替代全部方法的无代价最优解。

% Role: selected secondary endpoint results and ranking limits. Evidence: E1w.
\subsection{局部频谱结构与频带ILD}
Hybrid在三项局部频谱端点上取得十方法最低均值。高频一阶差分、二阶差分和多尺度凹口深度误差分别为0.3997\,dB/bin、0.2146\,dB/bin$^2$和0.4877\,dB（表\ref{tab:spectral-results}）。它们反映相邻频点变化与局部凹口深度的重建情况，和高频绝对幅度误差不是同一个评价问题。因此，Hybrid在这些指标上均值领先，与FSP-AE在主要对侧高频误差上占优并不矛盾。

% E1w: secondary test means; no Hybrid-centered secondary significance test.
\begin{table}[htbp]
\centering\small
\setlength{\tabcolsep}{4pt}
\bicaption[tab:spectral-results]{表}{局部频谱与频带ILD的补充比较（44名测试被试，Q26）}{Tab.}{Supplementary local-spectral and band-ILD comparison (44 test subjects, Q26)}
\begin{tabular}{lrrrr}
\toprule
方法 & 高频一阶差分 & 高频二阶差分 & 凹口深度 & 频带ILD\\
 & dB/bin$\downarrow$ & dB/bin$^2\downarrow$ & dB$\downarrow$ & dB$\downarrow$\\
\midrule
SH only & 0.6135 & 0.2990 & 0.8679 & 4.5959 \\
SUpDEq SH & 0.6163 & 0.3230 & 0.8034 & 2.9904 \\
SUpDEq NN & 0.5769 & 0.3106 & 0.7399 & 2.5197 \\
SUpDEq Bary. & 0.5765 & 0.3092 & 0.7411 & 2.5630 \\
MCA & 0.6033 & 0.3231 & 0.7778 & 2.0562 \\
MCAR v3.5.1 & 0.4251 & 0.2501 & 0.4977 & 1.6231 \\
FSP-AE & 0.4050 & 0.2153 & 0.4987 & 1.6994 \\
RANF & 0.4051 & 0.2172 & 0.4893 & 1.7360 \\
Hybrid & \textbf{0.3997} & \textbf{0.2146} & \textbf{0.4877} & 1.6205 \\
Bounded E25 & 0.4231 & 0.2365 & 0.5056 & \textbf{1.5734} \\
\bottomrule
\end{tabular}
\par\smallskip
\parbox{\linewidth}{\footnotesize 均为被试均值，加粗仅表示最低均值。高频差分与凹口深度使用767个插值方向，频带ILD使用72个水平面插值方向；本表为补充描述，不提升为独立确认性证据。}
\end{table}


% Role: required secondary emphasis without equivalence inflation. Evidence: E1w.
全域ERB与频带ILD可作为另外两个重点观察项，但都不是最优均值。Hybrid的频带ILD为1.6205\,dB，次于Bounded的1.5734\,dB；与全域ERB一起，两者相对各自最优均值分别高2.95\%与3.00\%，不能按1\%容差合并为最佳。频带ILD的相对位置也不能抵消严格水平面ILD的代价。三项局部指标仍属于补充描述，现有结果未提供本文所需的直接Hybrid中心次要端点配对区间，不宣称其对每个方法都显著更优。主凹口位置、检出情况或感知定位也不能由这些均值推出。

% Role: component evidence separated from test and isolated ablation. Evidence: E4a/E4p.
\subsection{频谱细化器的验证集表现}
在同一44名验证被试及767方向上，Hybrid相对FiLM E130集成的三项主要频谱误差均降低（表\ref{tab:component-results}）。全域ERB、对侧$25^\circ$ ERB和对侧高频分别降低0.023122、0.016183和0.200758\,dB，配对区间均低于零，对应改善人数为44、43和44人。水平面ILD差为$-0.000787$\,dB，区间跨零。这与追加频谱细化后的谱误差改善一致，但不说明已消除测试集中的ILD代价。

% E4a/E4p: matched three-member ensembles on validation, not test.
\begin{table}[htbp]
\centering\small
\setlength{\tabcolsep}{3pt}
\bicaption[tab:component-results]{表}{FiLM与Hybrid的验证集比较（44名被试，767方向；单位dB）}{Tab.}{Validation comparison of FiLM and Hybrid (44 subjects, 767 directions; dB)}
\begin{tabular}{lrrrr}
\toprule
指标$\downarrow$ & FiLM E130 & Hybrid E190 & 均值差 & 95\%配对区间\\
\midrule
全域ERB & 0.8275 & 0.8044 & $-0.023122$ & $[-0.025642,-0.020640]$ \\
对侧$25^\circ$ ERB & 1.2308 & 1.2146 & $-0.016183$ & $[-0.019667,-0.012974]$ \\
对侧高频 & 3.7028 & 3.5020 & $-0.200758$ & $[-0.216183,-0.185650]$ \\
水平面ILD & 0.6306 & 0.6298 & $-0.000787$ & $[-0.002977,+0.001443]$ \\
\bottomrule
\end{tabular}
\par\smallskip
\parbox{\linewidth}{\footnotesize 两种模型均为三个成员的等权集成。差值为Hybrid$-$FiLM；区间采用既有10000次被试配对bootstrap、种子20260828，未校正多重比较。}
\end{table}


两侧都是三成员集成，差异不仅包含新增CNN结构，还包含该模块的额外训练与复合目标优化。该比较属于开发证据，不是对CNN结构、各损失项或集成效应的逐项因果隔离，也不与测试集均值混合报告。

% Role: input sensitivity with a common validation evaluation domain. Evidence: E8a/E8p.
\subsection{观测方向数量敏感性}
冻结的Q26模型对观测数量变化并不表现为单调收益。在共同743个评价方向上，Q14相对Q26的四项误差均增加，配对区间均高于零；前三项频谱指标在44名被试上全部退化（表\ref{tab:q-results}）。Q50的全域与对侧$25^\circ$ ERB分别增加0.016283和0.051087\,dB，区间分别为$[+0.008783,+0.023196]$和$[+0.038154,+0.064047]$\,dB。高频与ILD的均值变化则分别为$+0.008602$和$+0.021770$\,dB，区间均含零。

% E8a: shared 743-direction validation domain, not the 767-direction domain.
\begin{table}[htbp]
\centering\small
\setlength{\tabcolsep}{4pt}
\bicaption[tab:q-results]{表}{冻结Hybrid的观测方向敏感性（44名验证被试，共同743方向）}{Tab.}{Input-direction sensitivity of frozen Hybrid (44 validation subjects, 743 shared directions)}
\begin{tabular}{lrrrr}
\toprule
输入 & 全域ERB & 对侧$25^\circ$ ERB & 对侧高频 & 水平面ILD\\
 & dB$\downarrow$ & dB$\downarrow$ & dB$\downarrow$ & dB$\downarrow$\\
\midrule
Q14 & $1.1057 \pm 0.1864$ & $1.6650 \pm 0.2485$ & $4.0345 \pm 0.3612$ & $0.7531 \pm 0.1896$ \\
Q26 & $0.8008 \pm 0.1614$ & $1.2123 \pm 0.1895$ & $3.4920 \pm 0.2761$ & $0.6232 \pm 0.1773$ \\
Q50 & $0.8171 \pm 0.1459$ & $1.2634 \pm 0.1895$ & $3.5006 \pm 0.2624$ & $0.6450 \pm 0.1527$ \\
\bottomrule
\end{tabular}
\par\smallskip
\parbox{\linewidth}{\footnotesize 数值为被试均值$\pm$样本标准差。统一排除全部Q50观测方向，网络参数不变；本表Q26不能与767方向表的均值直接比较。}
\end{table}


这些结果仅刻画Q26训练并冻结后的模型对不同输入集的响应，不证明更多测量方向一般有害，也不说明重新训练的Q50模型会产生同样变化。该表的Q26全域ERB为0.8008\,dB，来自验证集共同743方向，与组件表的0.8044\,dB及测试表的0.8175\,dB具有不同统计范围。

% Role: individual heterogeneity and retained failure case. Evidence: E2/E1s.
\subsection{个体差异与失败案例}
均值变化并不能替代个体表现。相对MCAR，Hybrid在全域ERB及对侧$25^\circ$ ERB均有39/44人改善，但前者总体配对区间仍含零；水平面ILD仅20/44人改善。测试被试P0339的Hybrid水平面ILD误差为6.0497\,dB，而同一被试的MCA和MCAR误差分别为0.7107与1.8838\,dB。该事后观察案例保留在原44人的统计中，不作剔除、修补或重测。它提示后续讨论须同时解释局部频谱收益和双耳线索风险，不能仅依据最优均值概括模型表现。
